\documentclass[10pt]{beamer}

% Språk och tecken
\usepackage[swedish]{babel}
\usepackage[utf8]{inputenc}
\usepackage[T1]{fontenc}

\usepackage{comment}

% Typsnitt – sans-serif för ren presentation
\usepackage{helvet}
\renewcommand{\familydefault}{\sfdefault}

% Tema: Använd default (enklare att anpassa färger)
\usetheme{default}
\usefonttheme{structurebold}

% FÄRGTEMA: Professionell grön AML/compliance-palett
% OBS! React-appen använder nu grön färgskala (brand-50 till brand-900)
% Definierad i tailwind.config.js:
%   - brand-50: #f7f9f8 (nästan vit med grön nyans)
%   - brand-100: #e8f0ed (mycket ljus grön)
%   - brand-600: #00704a (primär grön för knappar)
%   - brand-700: #005c3d (mörkare grön för hover)
%   - brand-800: #004d32 (mörk skogsgrön för headers)
%   - brand-900: #1a3a2e (mörkaste för text)
% 
% Detta LaTeX-dokument använder fortfarande varm beige/terracotta-palett
% för presentationsformat. Vid implementation i React används den gröna paletten.

\definecolor{warmbeige}{RGB}{245, 240, 230}     % Ljus bakgrund
\definecolor{warmbrown}{RGB}{101, 73, 50}        % Mörk text & accent
\definecolor{terracotta}{RGB}{204, 102, 85}      % För knappar / viktig info
\definecolor{olive}{RGB}{120, 110, 90}           % Sekundär accent

% Bakgrund och text
\setbeamercolor{background canvas}{bg=warmbeige}
\setbeamercolor{normal text}{fg=warmbrown}
\setbeamercolor{structure}{fg=warmbrown}

% Titel och frame-titel
\setbeamercolor{title}{fg=white, bg=warmbrown}
\setbeamercolor{frametitle}{fg=white, bg=warmbrown}

% Knappar och länkar
\setbeamercolor{button}{bg=terracotta, fg=white}
\setbeamercolor{button border}{use=button, fg=button.bg}

% Block (t.ex. exempel, varningar)
\setbeamercolor{block title}{fg=white, bg=olive}
\setbeamercolor{block body}{bg=warmbeige!70!white}

% Stil för knappar
\setbeamertemplate{buttons}[rounded]

% Typsnittsstorlek
\setbeamerfont{frametitle}{size=\Large}
\setbeamerfont{title}{size=\LARGE}

% Titel
\title[Kundonboarding]{Kundonboarding: Accept eller Reject}
\author[Celestial AB]{Celestial AB}
\institute{Risk \& Compliance}
\date{\today}

% Snygga länkar utan röda ramar
\hypersetup{
  colorlinks=true,
  linkcolor=terracotta,
  urlcolor=terracotta
}

\begin{document}

% ========================================================================
% DEL 1: AUTENTISERINGSSIDOR (UTAN SIDEBAR)
% ========================================================================
% Dessa sidor visas INNAN användaren är inloggad och har INGEN sidebar.
% De hanterar registrering, verifiering, inloggning och lösenordsåterställning.
% ========================================================================

% Hero-sektionen (förbättrad, mer lättläst och säljande)
\begin{frame}[plain]
  % Knappar uppe till höger
  \begin{flushright}
    \vspace{-0.8cm}
    \hyperlink{login}{\beamerbutton{Logga in}} \hspace{0.3cm}
    \hyperlink{register}{\beamerbutton{Registrera}}
  \end{flushright}

  \vspace{1.2cm}
  \centering
  {\usebeamerfont{title}\color{warmbrown}\textbf{Onboardingstöd för redovisnings- och revisionsbyråer}}

  \vspace{0.8cm}

  % Kort introduktion / värde
  \begin{block}{Snabb översikt}
    Automatisera riskbedömningen och säkerställ att onboarding av nya klienter sker enligt penningtvättslagen och tillsynsmyndighetens riktlinjer.
  \end{block}

  \vspace{0.5cm}

  % AI och automatisering
  \begin{block}{AI-stöd och kontroll}
    Vår AI-baserade programvara samlar automatiskt all nödvändig information, kontrollerar klienten mot penningtvättsregister och företagsbankkonton för att upptäcka \textit{layering}, och hjälper dig att skriva och signera avtal med BankID.
  \end{block}

  \vspace{0.5cm}

  % Compliance och lag
  \begin{block}{Compliance och lagkrav}
    Följ vår mall så sker onboardingprocessen enligt \textit{Länsstyrelsen Stockholms författningssamling 01FS 2024–20}.
  \end{block}

  \vspace{0.5cm}

  % Funktioner (kortare, visuellt uppdelat)
  \begin{block}{Vad du får hjälp med}
    \begin{columns}[T]
      \begin{column}{0.48\textwidth}
        \textbullet\ Ägarstruktur \\
        \textbullet\ Alternativ verklig huvudman \\
        \textbullet\ Bankkonton \\
        \textbullet\ Finansiell information \\
        \textbullet\ Folkbokföringsregistret \\
        \textbullet\ Företagsdokument
      \end{column}
      \begin{column}{0.48\textwidth}
        \textbullet\ Aktiefakta \\
        \textbullet\ Arbetsställen \\
        \textbullet\ Fastighetsinformation \\
        \textbullet\ Firmatecknare \\
        \textbullet\ Företagsärenden \\
        \textbullet\ Koncernträd
      \end{column}
    \end{columns}
  \end{block}

  \vspace{0.5cm}

  % CTA
  \beamerbutton{Börja onboarding nu} \hspace{0.5cm} \beamerbutton{Se demo}
\end{frame}

% Herosektion fortsättning – mer lättläst
\begin{frame}[plain]
  \frametitle{Onboardingstöd – fortsättning}
  \footnotesize
  \begin{minipage}{0.92\textwidth}
    \textbf{Få full kontroll över klienten med våra automatiserade kontroller:}
    \medskip
    \begin{columns}[T]
      \begin{column}{0.48\textwidth}
        \textbullet\ Kreditbeslut \\
        \textbullet\ Penningtvättsregistret \\
        \textbullet\ Rating av företag \\
        \textbullet\ Sanktionslistor \\
        \textbullet\ Styrelseuppdrag \\
        \textbullet\ Verksamhetsbeskrivning
      \end{column}
      \begin{column}{0.48\textwidth}
        \textbullet\ Näringsförbud \\
        \textbullet\ PEP \\
        \textbullet\ Riskindikatorer \\
        \textbullet\ Styrelsemedlemmar \\
        \textbullet\ Verklig huvudman
      \end{column}
    \end{columns}
  \end{minipage}
\end{frame}

%Registreringssidan

\begin{frame}[label=register]
  \frametitle{Skapa konto}

  \vspace{0.5cm}
  \begin{minipage}{0.9\textwidth}
    \small
    \textbf{E-post:} \underline{\hspace{6.5cm}} \\
    \textbf{Lösenord:} \underline{\hspace{6.5cm}} \\
    \footnotesize\textit{(Minst 12 tecken, en siffra, en versal)} \\
    \textbf{Bekräfta lösenord:} \underline{\hspace{6.5cm}} \\

    \vspace{0.8cm}
    $\Box$ Jag bekräftar att jag är en människa (inte en bot).
    \footnotesize
    \textit{(Vi använder automatiska skydd, t.ex. Cloudflare, för att förhindra missbruk och spam.)}

    \vspace{1.2cm}
    \centering
    \beamerbutton{Skapa konto}
    
    \vspace{1.2cm}
    \textcolor{warmbrown}{\textbf{eller}}
    
    \vspace{0.6cm}
    \colorbox{white}{%
      \parbox{5cm}{\centering\footnotesize\textbf{Logga in med Google}}%
    }
  \end{minipage}
  
  \vspace{0.8cm}
  \begin{center}
    \footnotesize
    Har du redan konto? \hyperlink{login}{\beamergotobutton{Logga in}}
  \end{center}
\end{frame}

%Sidan som simulerar registreringskoder erhållna via e-post
\begin{frame}[label=verify]
  \frametitle{Verifiera din registrering}

  \vspace{0.5cm}
  \begin{minipage}{0.9\textwidth}
    \small
    \textbf{Ange den registreringskod du fått via e-post eller sms:} \underline{\hspace{5cm}} \\

    \vspace{0.8cm}
    \begin{block}{Info}
      Av säkerhetsskäl och på grund av problem med e-postlänkar (t.ex. Sendgrid) används nu registreringskoder istället för klickbara länkar.
    \end{block}

    \vspace{1.2cm}
    \centering
    \beamerbutton{Verifiera kod}
    
    \vspace{1.2cm}
    \textcolor{warmbrown}{\textbf{eller}}
    
    \vspace{0.6cm}
    \colorbox{white}{%
      \parbox{5cm}{\centering\footnotesize\textbf{Skicka ny kod}}%
    }
  \end{minipage}
\end{frame}

%Login sidan

\begin{frame}[label=login]
  \frametitle{Logga in}

  \vspace{0.5cm}
  \begin{minipage}{0.9\textwidth}
    \small
    \textbf{E-post:} \underline{\hspace{6.5cm}} \\
    \textbf{Lösenord:} \underline{\hspace{6.5cm}} \\
    
    \vspace{0.3cm}
    \footnotesize
    \hyperlink{forgotpassword}{\textcolor{terracotta}{Glömt lösenord?}}

    \vspace{0.8cm}
    \centering
    \beamerbutton{Logga in}

    \vspace{1.2cm}
    \textcolor{warmbrown}{\textbf{eller}}

    \vspace{0.6cm}
    \colorbox{white}{%
      \parbox{5cm}{\centering\footnotesize\textbf{Logga in med Google}}%
    }
  \end{minipage}

  \vspace{0.8cm}
  \begin{center}
    \footnotesize
    Har du inget konto? \hyperlink{register}{\beamergotobutton{Registrera}}
  \end{center}
\end{frame}

% Sida Glömt lösenord (ForgotPasswordSlide i React, route: /forgot-password)
\begin{frame}[label=forgotpassword]
  \frametitle{Återställ lösenord}

  \vspace{0.5cm}
  \begin{minipage}{0.9\textwidth}
    \small
    \textbf{E-post:} \underline{\hspace{6.5cm}} \\

    \vspace{1cm}
    \centering
    \beamerbutton{Skicka återställningskod}

    \vspace{1cm}
    \footnotesize
    En 6-siffrig kod kommer skickas till din e-post. Koden är giltig i 15 minuter.
  \end{minipage}

  \vspace{0.8cm}
  \begin{center}
    \footnotesize
    \hyperlink{login}{\beamergotobutton{Tillbaka till login}}
  \end{center}
\end{frame}

% Sida: Ange nytt lösenord (ResetPasswordSlide i React, route: /reset-password)
\begin{frame}[label=resetpassword]
  \frametitle{Ange nytt lösenord}

  \vspace{0.5cm}
  \begin{minipage}{0.9\textwidth}
    \small
    \textbf{Verifieringskod (6 siffror):} \underline{\hspace{6.5cm}} \\
    \vspace{0.5cm}
    \textbf{Nytt lösenord:} \underline{\hspace{6.5cm}} \\
    \footnotesize\textit{(Minst 12 tecken, en siffra, en versal)} \\
    \vspace{0.3cm}
    \textbf{Bekräfta lösenord:} \underline{\hspace{6.5cm}} \\

    \vspace{1cm}
    \centering
    \beamerbutton{Återställ lösenord}

    \vspace{1cm}
    \footnotesize
    Fick du inte koden? \hyperlink{forgotpassword}{\textcolor{terracotta}{Skicka ny kod}}
  \end{minipage}
\end{frame}

% ========================================================================
% DEL 2: ONBOARDING SLIDES (MED SIDEBAR)
% ========================================================================
% Dessa slides visas EFTER inloggning och har sidebar med progressindikator.
% De utgör själva onboarding-processen där företagsdata samlas in steg för steg.
% ========================================================================

% --- Kommentar: Inloggad vy med vänster sidebar och ikoner ---
% När användaren är inloggad visas följande layout:
%
% 1. Vänster sidebar
%    - Expanderbar
%    - Mini-thumbnails över slidsen / steg i onboarding
%    - Menylista med kugghjul för inställningar
%
% 2. Ikoner uppe i vänster sidebar eller över content-ytan
%    - LLM: Klick på denna fäller ut en chattpanel till höger.
%           LLM får via ett JSON-paket all data som appen hittills hämtat
%           för att kunna svara på frågor efter bästa förmåga.
%
%    - Dokumentation: Klick fäller ut dokumentationspanel där användaren
%           ser hur programmet används och vilka tester som körs.
%
%    - Support: Klick fäller ut en supportpanel med möjlighet till
%           chatt och delning av skärm (likt Fortnox). 
%           Knappen i den utfällda menyn heter "Supportinloggning".
%           Symbol: ring med frågetecken.
%
% 3. Content-ytan
%    - Visar huvudfunktioner, formulär och PDF/graph-visualisering
%    - Eventuellt LLM / supportpanel fälls ut till höger, påverkar
%      content-bredd dynamiskt.
%
% Kommentar: Denna design följer principen om modulär sidebar och
% dynamiska paneler, vilket möjliggör både onboarding och interaktiv
% assistans utan att navigera bort från huvudytan.

% Slide 1: Inledning och bakgrund (IntroSlide i React, route: /inledning) - FÖRSTA SLIDE MED SIDEBAR
% Före detta finns Auth-slides (0-3): Hero (/, HeroSlide), Login (/login), Register (/register), Verify (/verify)
\begin{frame}[label=intro]
  \frametitle{Inledning och bakgrund}

  \small
  Denna onboarding säkerställer att byrån uppfyller penningtvättslagstiftningens krav vid antagandet av nya kunder. Processen efterlever tillsynsmyndighetens krav där verksamhetsutövaren (byrån) är skyldig att redovisa hur man säkerställt att byrån inte gör sig skyldig till penningtvätt.\\
  Länsstyrelserna, som ansvarar för tillsynen, har skärpt kraven på redovisningsbyråer och utfärdar sanktionsavgifter på hundratusentals kronor vid bristande efterlevnad.\\
  Det är dessa myndighetskrav som tvingar oss att ställa specifika frågor och att spara dokumentationen i minst fem år.\\
  Processen är därmed en följd av myndighetskrav och syftar till att förebygga penningtvätt och ekonomisk brottslighet.

  \vspace{1.2cm}
  \begin{flushright}
    \hyperlink{nextslide}{\beamerbutton{Nästa}}
  \end{flushright}
\end{frame}

% Slide 2: Riskfrågor - Uppdragets syfte och art (RiskfragorSlide i React, route: /riskfragor)
\begin{frame}[label=uppdrag]
  \frametitle{Uppdragets syfte och art}

  \small
  Vad exakt är det kunden vill ha hjälp med?

  \vspace{0.5cm}
  \textbf{Typ av uppdrag:}
  \begin{itemize}
    \item[$\Box$] Bokföring
    \item[$\Box$] Lönehantering
    \item[$\Box$] Rådgivning
    \item[$\Box$] Deklaration
    \item[$\Box$] Annat
  \end{itemize}

  \vspace{0.3cm}
  \textbf{Frekvens:}
  \begin{itemize}
    \item[$\Box$] Löpande
    \item[$\Box$] Engångsuppdrag
    \item[$\Box$] Periodiskt
  \end{itemize}

  \vspace{0.3cm}
  \textbf{Eventuella särskilda omständigheter:}
  \begin{block}{Textfält}
    \underline{\hspace{10cm}}
  \end{block}

  \vspace{1cm}
  \begin{flushright}
    \hyperlink{verksamhet}{\beamerbutton{Nästa}}
  \end{flushright}
\end{frame}


% Slide 2 (fortsättning): Kundens verksamhetsart - ingår i RiskfragorSlide
\begin{frame}[label=verksamhet]
  \frametitle{Kundens verksamhetsart}

  \small
  Penningtvättslagen (PTL) kräver att redovisningsbyråer fastställer varje kunds riskprofil.\\
  Tillsynsmyndigheten Länsstyrelsen Stockholm ställer tydliga krav på att byrån ska kunna visa att man förstått kundens affärsidé, verksamhetshistorik samt kund- och leverantörsstruktur.\\
  Det är viktigt att kunna bevisa detta vid en eventuell granskning. Om byrån inte kan styrka att man tydligt övertygat sig om att kundens affärer är rena, kan stora sanktionsavgifter tillfogas – även om kunden själv är helt felfri. Detta kallas "klandervärt risktagande" och innebär ansvar för byrån utan att brottspengar behöver vara inblandade.\\
  Se gärna mer information och exempel på faktiska sanktionsbeslut hos 
  \href{https://www.lansstyrelsen.se/stockholm/samhalle/betalning-ekonomi-och-pengar/forhindra-penningtvatt-och-finansiering-av-terrorism/ingripanden-och-sanktioner-penningtvatt.html}{Länsstyrelsen Stockholm}.

  \vspace{1cm}
  \begin{block}{Syfte}
    Vi ställer dessa frågor för att följa lagkrav och skydda både dig och byrån – inte av nyfikenhet.
  \end{block}

  \vspace{0.8cm}
  \begin{flushright}
    \hyperlink{nextslide}{\beamerbutton{Nästa}}
  \end{flushright}
\end{frame}

% Slide 2 (fortsättning): Frågor som stödjer riskbedömning - ingår i RiskfragorSlide
\begin{frame}[label=riskfragor]
  \frametitle{Frågor som stödjer riskbedömning}

  \small
  Flera av dessa frågor har lagstöd och hjälper oss att bedöma risken:
  \begin{itemize}
    \item Vad är företagets huvudsakliga affärsidé?
    \item Vilka typer av kunder har företaget? (Privatpersoner / Företag / Offentlig sektor)
    \item Har företaget återkommande utländska affärspartners?
    \item Vilka är de största leverantörerna, och var är de etablerade?
    \item Har verksamheten ändrats på senare tid?
    \item Är du eller någon i företaget en PEP (person i politiskt utsatt ställning)?
  \end{itemize}

  \vspace{0.5cm}
  	\textbf{Organisationsnummer:} \underline{\hspace{5cm}}
  \footnotesize
  \begin{block}{Info}
    Organisationsnumret används för att hämta officiell information från Bolagsverket eller Roaring.io.
  \end{block}

  \vspace{0.5cm}
  	\textbf{Personnummer:} \underline{\hspace{5cm}}
  \footnotesize
  \begin{block}{Info}
    Personnumret används för att hämta officiell information från Bolagsverket eller Roaring.io.
  \end{block}

  \vspace{0.8cm}
  \begin{flushright}
    \hyperlink{nextslide}{\beamerbutton{Nästa}}
  \end{flushright}
\end{frame}



% Slide 3: Identitetskontroll och dokumentation (IdentitetskontrollSlide i React, route: /identitet)
% Denna slide är alltid med och följer 3 kap. 2 och 4 §§.
\begin{frame}[label=identitetskontroll]
  \frametitle{Identitetskontroll och dokumentation}

  \small
  För att uppfylla penningtvättslagens krav måste identiteten kontrolleras och dokumenteras enligt följande:
  \begin{itemize}
    \item \textbf{Fysisk person:} Pass, körkort eller annan identitetshandling med foto utfärdad av myndighet, tillförlitlig elektronisk legitimation, eller dokument/uppgifter från oberoende och tillförlitliga källor.
    \item \textbf{Juridisk person:} Registreringsbevis eller motsvarande utdrag (ej äldre än en vecka), samt uppgifter om företrädare.
    \item \textbf{Ombud/företrädare:} Fullmakt eller motsvarande behörighetshandling.
  \end{itemize}

  \vspace{0.5cm}
  \textbf{Dokumentation:}
  \begin{itemize}
    \item Anteckna identitetshandlingens nummer och giltighetstid eller bevara kopia.
    \item Bevara kopia av elektronisk legitimation eller andra dokument som legat till grund för kontrollen.
    \item Kopia av registreringsbevis, fullmakt och andra relevanta handlingar ska sparas.
    \item All dokumentation ska sparas minst fem år och vara sökbar.
  \end{itemize}

  \vspace{0.8cm}
  \begin{minipage}{0.45\textwidth}
    \begin{flushleft}
      % Simulerad knapp för att ta foto med webbkamera
      \beamerbutton{Ta foto med webbkamera}
      \footnotesize
      	\textit{(Personen håller upp sitt ID/körkort)}
    \end{flushleft}
  \end{minipage}%
  \hfill
  \begin{minipage}{0.45\textwidth}
    \begin{flushright}
      \hyperlink{nextslide}{\beamerbutton{Nästa}}
    \end{flushright}
  \end{minipage}
\end{frame}

% Slide 4: Identitetskontroll – sammanfattande tabell (KontrolltabellSlide i React, route: /kontrolltabell)
% Tabell enligt Länsstyrelsen Stockholms författningssamling 01FS 2024:20 och penningtvättslagen (2017:630)
\begin{frame}[label=kontrolltabell]
  \frametitle{Identitetskontroll – sammanfattande tabell}
  \footnotesize
  \textit{Tabellen nedan sammanfattar kraven enligt Länsstyrelsen Stockholms författningssamling 01FS 2024:20 och lagen (2017:630) om åtgärder mot penningtvätt och finansiering av terrorism.}
  \vspace{0.3cm}
  \begin{tabular}{|p{2.8cm}|p{5.2cm}|p{5.2cm}|}
    \hline
    \textbf{Vem kontrolleras} & \textbf{Vad kontrolleras} & \textbf{Dokumentation av kontrollen} \\
    \hline
    Fysisk person & 
    1. Pass, körkort eller annan identitetshandling med foto utfärdad av myndighet eller tillförlitlig utfärdare.\\
    2. Tillförlitlig elektronisk legitimation.\\
    3. Andra dokument/uppgifter från oberoende och tillförlitliga källor. Vid svårigheter, använd flera källor. &
    1. Anteckna identitetshandlingens nummer och giltighetstid eller bevara kopia.\\
    2. Bevara kopia av bekräftelsen på elektronisk legitimation.\\
    3. Bevara kopia av dokument/uppgifter som legat till grund för kontrollen. \\
    \hline
    Fysisk person på distans &
    1. Tillförlitlig elektronisk legitimation.\\
    2. Namn, personnummer/samordningsnummer och adress mot externa register/intyg/oberoende källor, samt\\
    a) Skicka bekräftelse till folkbokföringsadress, eller\\
    b) Inhämta vidimerad kopia av identitetshandling. &
    1. Bevara kopia av bekräftelsen på elektronisk legitimation.\\
    2. Bevara kopia av dokument/uppgifter som legat till grund för kontrollen samt kopia av bekräftelsebrev eller vidimerad kopia. \\
    \hline
    Företrädare, ombud eller motsvarande &
    1. Se identitetskontroll för fysisk person eller fysisk person på distans.\\
    2. Fullmakt, förordnande eller motsvarande behörighetshandling. &
    1. Se dokumentationskrav för fysisk person eller fysisk person på distans.\\
    2. Bevara kopia av fullmakt, förordnande eller motsvarande behörighetshandling. \\
    \hline
    Juridisk person &
    Registreringsbevis, registerutdrag eller uppgifter från andra tillförlitliga och oberoende källor. &
    Bevara kopia av de kontrollerade handlingarna. \\
    \hline
    Verklig huvudman &
    Se identitetskontroll för fysisk person eller fysisk person på distans. &
    Se dokumentationskrav för fysisk person eller fysisk person på distans. \\
    \hline
  \end{tabular}
  \vspace{0.3cm}
  \footnotesize
  \textit{Se även:}
  \begin{itemize}
    \item Länsstyrelsen Stockholms författningssamling 01FS 2024:20, 3 kap. 4 § och tabell 1
    \item Lagen (2017:630) om åtgärder mot penningtvätt och finansiering av terrorism
  \end{itemize}
\end{frame}



  % Slide 5: Skärpt kundkännedom vid PEP (PEPSlide i React, route: /pepfordjupning)
  % OBS! Denna slide visas dynamiskt om kunden eller någon i företaget är PEP.
  \begin{frame}[label=pepfordjupning]
    \frametitle{Skärpt kundkännedom – PEP}

    \small
    Eftersom du eller någon i företaget är en PEP (person i politiskt utsatt ställning) krävs fördjupad kontroll enligt penningtvättslagen. Vänligen besvara följande frågor:
    \begin{itemize}
      \item Vad är ursprunget till de medel som ska användas i affärsförbindelsen?
      \item Beskriv affärsverksamhetens syfte och art.
      \item Vilka ekonomiska medel och resurser har företaget tillgång till?
      \item Finns det ytterligare dokumentation som styrker medlens ursprung?
      \item Har du eller någon i företaget affärsrelationer med högriskländer?
      \item Finns det andra omständigheter som kan påverka riskbedömningen?
    \end{itemize}

    \vspace{0.8cm}
    \begin{flushright}
      \hyperlink{nextslide}{\beamerbutton{Nästa}}
    \end{flushright}
  \end{frame}

% Slides 6 till 10: Jämförelse mot Bolagsverket/Roaring.io

% Slide 6: Verksamhetsbeskrivning & verksamhetskoder (VerksamhetSlide i React, route: /verksamhet) - RESULTATSLIDE från Roaring API
\begin{frame}[label=verksamhetskod]
  \frametitle{Verksamhetsbeskrivning och SNI-koder}
  \small
  \textbf{Beskrivning:} \underline{\hspace{8cm}} \\
  \textbf{Bransch/SNI-kod:} \underline{\hspace{8cm}} \\
  \textbf{Sekundära namn:} \underline{\hspace{8cm}} \\
  \vspace{0.5cm}
  \begin{block}{Jämförelse}
    Om beskrivningen eller bransch inte stämmer med kundens uppgifter, diskutera och dokumentera avvikelsen.
  \end{block}
  \vspace{0.8cm}
  \begin{flushright}
    \hyperlink{nextslide}{\beamerbutton{Nästa}}
  \end{flushright}
\end{frame}

% Slide 7: Ägarstruktur & verklig huvudman (AgarstrukturSlide i React, route: /agarstruktur) - RESULTATSLIDE från Roaring API
\begin{frame}[label=agarstruktur]
  \frametitle{Ägarstruktur och verklig huvudman}
  \small
  \textbf{Ägare:} \underline{\hspace{8cm}} \\
  \textbf{Verklig huvudman:} \underline{\hspace{8cm}} \\
  \textbf{Alternativa huvudmän:} \underline{\hspace{8cm}} \\
  \textbf{Ägarandelar:} \underline{\hspace{8cm}} \\
  \vspace{0.5cm}
  \begin{block}{Jämförelse}
    Om ägaruppgifter avviker från kundens svar, diskutera och dokumentera.
  \end{block}
  \vspace{0.8cm}
  \begin{flushright}
    \hyperlink{nextslide}{\beamerbutton{Nästa}}
  \end{flushright}
\end{frame}

% Slide 8: Styrelsemedlemmar & firmatecknare (StyrelseSlide i React, route: /styrelse) - RESULTATSLIDE från Roaring API
\begin{frame}[label=styrelse]
  \frametitle{Styrelsemedlemmar och firmatecknare}
  \small
  \textbf{Styrelse:} \underline{\hspace{8cm}} \\
  \textbf{VD:} \underline{\hspace{8cm}} \\
  \textbf{Firmatecknare:} \underline{\hspace{8cm}} \\
  \textbf{Signaturrätt:} \underline{\hspace{8cm}} \\
  \textbf{Historik:} \underline{\hspace{8cm}} \\
  \vspace{0.5cm}
  \begin{block}{Jämförelse}
    Om styrelseuppgifter avviker från kundens svar, diskutera och dokumentera.
  \end{block}
  \vspace{0.8cm}
  \begin{flushright}
    \hyperlink{nextslide}{\beamerbutton{Nästa}}
  \end{flushright}
\end{frame}

% Slide 9: Riskindikatorer & sanktionslistor (RiskindikatorerSlide i React, route: /riskindikatorer) - RESULTATSLIDE från Roaring API
\begin{frame}[label=riskindikator]
  \frametitle{Riskindikatorer och sanktionslistor}
  \small
  \textbf{Riskpoäng:} \underline{\hspace{8cm}} \\
  \textbf{Larm:} \underline{\hspace{8cm}} \\
  \textbf{Sanktionslistor:} \underline{\hspace{8cm}} \\
  \textbf{PEP:} \underline{\hspace{8cm}} \\
  \vspace{0.5cm}
  \begin{block}{Jämförelse}
    Om riskindikatorer eller sanktionslistor avviker från kundens svar, diskutera och dokumentera.
  \end{block}
  \vspace{0.8cm}
  \begin{flushright}
    \hyperlink{nextslide}{\beamerbutton{Nästa}}\begin{frame}[label=login]
  \frametitle{Logga in}
  
  ...
  \textbf{Lösenord:} \underline{\hspace{6.5cm}} \\
  
  \vspace{0.3cm}
  \footnotesize
  \hyperlink{forgotpassword}{\textcolor{terracotta}{Glömt lösenord?}}
  
  \vspace{0.8cm}
  \beamerbutton{Logga in}
  ...
\end{frame}

\begin{frame}[label=forgotpassword]
  \frametitle{Återställ lösenord}
  
  \textbf{E-post:} \underline{\hspace{6.5cm}} \\
  
  \vspace{0.8cm}
  \beamerbutton{Skicka återställningskod}
  
  \vspace{1cm}
  \footnotesize
  En 6-siffrig kod kommer skickas till din e-post.
\end{frame}

\begin{frame}[label=resetpassword]
  \frametitle{Ange ny lösenord}
  
  \textbf{Verifieringskod:} \underline{\hspace{6.5cm}} \\
  \textbf{Nytt lösenord:} \underline{\hspace{6.5cm}} \\
  \textbf{Bekräfta lösenord:} \underline{\hspace{6.5cm}} \\
  
  \vspace{0.8cm}
  \beamerbutton{Återställ lösenord}
\end{frame}
  \end{flushright}
\end{frame}

% Slide 10: Övriga datapunkter (OvrigaDataSlide i React, route: /ovrigadata) - RESULTATSLIDE från Roaring API
\begin{frame}[label=ovrigt]
  \frametitle{Övriga datapunkter}
  \small
  \textbf{Fastigheter:} \underline{\hspace{8cm}} \\
  \textbf{Engagemang:} \underline{\hspace{8cm}} \\
  \textbf{Rättsliga ärenden:} \underline{\hspace{8cm}} \\
  \vspace{0.5cm}
  \begin{block}{Jämförelse}
    Om övriga datapunkter avviker från kundens svar, diskutera och dokumentera.
  \end{block}
  \vspace{0.8cm}
  \begin{flushright}
    \hyperlink{nextslide}{\beamerbutton{Nästa}}
  \end{flushright}
\end{frame}

% =======================================================================
% FLÖDE: Bokföringsdata → Ekonomisk rådgivning (värdeskapande) → Djupgranskning
% =======================================================================
% OBS: Slide 13 kombinerad inmatning av IBAN + SIE-fil (Fortnox/Visma/Bokio + manuell uppladdning)
% Detta motsvarar BokforingDataSlide i React (slide 14, route: /bokforing)
% React-implementering har ytterligare slides före detta:
% - Slides 0-3: Auth (Hero, Login, Register, Verify)
% - Slide 4: Inledning (IntroSlide, route: /inledning)
% - Slides 5-8: Riskfrågor (RiskfragorSlide, IdentitetskontrollSlide, KontrolltabellSlide, PEPSlide)
% - Slides 9-13: Result slides (Verksamhet, Ägarstruktur, Styrelse, Riskindikatorer, Övriga data)
%
% EKONOMISK RÅDGIVNING (Slides 11-14): Värdeskapande analys FÖRE potentiell avslag
% Slide 14: Likviditetsanalys – Graf över likviditet från bankkonto (LikviditetsanalysSlide)
% Slide 15: Omsättningsanalys – Omsättning per räkenskapsår från SIE (OmsattningsanalysSlide)
% Slide 16: Resultatanalys – Årets resultat per räkenskapsår från SIE (ResultatanalysSlide)
% Slide 17: Branschjämförelse – Företagets nyckeltal vs SCB-data för branschen (BranschjamforelseSlide)
%
% DJUPGRANSKNING (Slides 19-23): Detaljerad kontroll som kan leda till avslag
% Slide 18: Bokföringsanalys – AI-granskning av verifikationer (BokforingsanalysSlide)
% Slide 19: Riskbedömning & Besked – Slutligt besked (RiskbedomningSlide)
% Slide 20: Kundens skyldigheter – Vad kunden måste uppfylla (SkyldigheterSlide)
% Slide 21: Avtalsvillkor och godkännande – BankID-signering (AvtalSlide)
% Slide 22: E-postbekräftelse och leverans av dokument (DocumentDeliverySlide)
%
% Config.json styr:
%   - Antal år för SIE-filhämtning (standard 7 år, kan minskas för nystartade företag)
%   - Vilka risktester som är aktiverade
%   - Textanpassningar och valfria fält
%
% Resultatslides (6-10) använder grönt tema i React-implementationen.
% =======================================================================

% Slide 11: Företagsdokumentation (ForetagsdokumentationSlide i React, route: /dokumentation)
\begin{frame}[label=foretagsdokumentation]
  \frametitle{Företagsdokumentation}
  \small
  För att verifiera företagets existens och styrelsemedlemmar behöver vi grundläggande företagsdokumentation.
  
  \vspace{0.3cm}
  \begin{block}{Varför behövs detta?}
    \footnotesize
    ✅ \textbf{Registreringsbevis} verifierar företagets existens hos Bolagsverket\\
    ✅ Ger oss \textbf{styrelsemedlemmar} och \textbf{VD} med personnummer (för närstående-kontroll)\\
    ✅ Bekräftar företagets \textbf{registrerade adress} (för leveransadress-validering)\\
    ✅ \textbf{Standard vid onboarding} enligt PTL 2 kap. 7 § (kundkännedom)
  \end{block}
  
  \vspace{0.3cm}
  \begin{minipage}{0.48\textwidth}
    \textbf{Registreringsbevis från Bolagsverket:}\\
    \vspace{0.2cm}
    
    % Dropyta för registreringsbevis (obligatoriskt)
    \fbox{\begin{minipage}{4.2cm}
      \centering
      \vspace{0.3cm}
      📄\\
      \vspace{0.2cm}
      \footnotesize\textbf{Dra och släpp PDF här}\\
      \vspace{0.1cm}
      \tiny{eller klicka för att välja fil}\\
      \vspace{0.3cm}
    \end{minipage}}\\
    \vspace{0.2cm}
    \footnotesize\textcolor{red}{* Obligatoriskt (max 3 mån gammalt)}
    
  \end{minipage}
  \hfill
  \begin{minipage}{0.48\textwidth}
    \textbf{Årsredovisning (valfritt):}\\
    \vspace{0.2cm}
    
    % Dropyta för årsredovisning (valfritt)
    \fbox{\begin{minipage}{4.2cm}
      \centering
      \vspace{0.3cm}
      📊\\
      \vspace{0.2cm}
      \footnotesize\textbf{Dra och släpp PDF här}\\
      \vspace{0.1cm}
      \tiny{eller klicka för att välja fil}\\
      \vspace{0.3cm}
    \end{minipage}}\\
    \vspace{0.2cm}
    \footnotesize\textcolor{gray}{Valfritt men rekommenderat}
  \end{minipage}
  
  \vspace{0.3cm}
  \begin{block}{Vad händer med dokumenten?}
    \footnotesize
    Vi \textbf{parsar automatiskt} registreringsbeviset med OCR/PDF-parsing för att extrahera:\\
    • Organisationsnummer → Verifieras mot Bolagsverket GRATIS API\\
    • Företagsnamn → Verifieras mot angivna uppgifter\\
    • Styrelseledamöter + personnummer → Används för närstående-kontroll (SPAR/Roaring.io)\\
    • Registrerad adress → Används för leveransadress-validering i bokföringsanalysen
  \end{block}
  
  \vspace{0.3cm}
  \begin{flushright}
    \hyperlink{bokforing_bank}{\beamerbutton{Nästa}}
  \end{flushright}
\end{frame}

% Slide 13: Bokföringsdata och bankkonto (BokforingDataSlide i React, route: /bokforing)
\begin{frame}[label=bokforing_bank]
  \frametitle{Bokföringsdata och Skattekonto}
  \small
  För att göra en fullständig riskbedömning behöver vi tillgång till företagets bokföringsunderlag och skattedata.
  
  \vspace{0.3cm}
  \begin{minipage}{0.48\textwidth}
    \textbf{Bokföringsdata (SIE-filer):}\\
    \vspace{0.2cm}
    
    % Integration mot bokföringsprogram
    \beamerbutton{Fortnox} \quad
    \beamerbutton{Visma} \quad
    \beamerbutton{Bokio}\\
    \footnotesize\textit{(OAuth-integration, hämtar automatiskt)}
    
    \vspace{0.2cm}
    \textbf{eller ladda upp manuellt:}\\
    \vspace{0.2cm}
    
    % Dropyta för manuell uppladdning
    \fbox{\rule{4.5cm}{1.2cm}}\\
    \footnotesize\textit{Dra och släpp SIE-filer här}
    
  \end{minipage}
  \hfill
  \begin{minipage}{0.48\textwidth}
    \textbf{Skattekonto \& Deklarationer:}\\
    \vspace{0.2cm}
    
    \beamerbutton{Hämta via Skatteverket}\\
    \footnotesize\textit{(BankID-autentisering, hämtar automatiskt)}
    
    \vspace{0.2cm}
    \textbf{eller ladda upp CSV manuellt:}\\
    \vspace{0.2cm}
    
    \fbox{\rule{4.5cm}{1.2cm}}\\
    \footnotesize\textit{Skattekonto CSV från skatteverket.se}
  \end{minipage}
  
  \vspace{0.3cm}
  \begin{block}{Skatteverket OAuth2-flöde (ACG med BankID)}
    \footnotesize
    \textbf{1. Användare klickar:} "Hämta via Skatteverket" → OAuth2-redirect till Skatteverkets inloggningssida\\
    \textbf{2. BankID-autentisering:} Användaren loggar in med BankID \textit{(Skatteverket betalar, ingen kostnad för appen)}\\
    \textbf{3. Godkännande:} Användaren godkänner åtkomst till \textbf{både} Skattekonto \textbf{och} Inkomstdeklarationer (7 år)\\
    \textbf{4. Redirect tillbaka:} Auktorisationskod → access token → hämtar data via API\\
    \textbf{Scope:} \texttt{skattekonto:read inkomstdeklaration:read} (kombinerade i ETT anrop)\\
    \textbf{Fallback:} Om OAuth2 misslyckas, visa handledning för manuell CSV-nedladdning
  \end{block}
  
  \vspace{0.3cm}
  \begin{flushright}
    \hyperlink{bokforingsunderlag}{\beamerbutton{Nästa}}
  \end{flushright}
\end{frame}

% Slide 12: Bokföringsunderlag (BokforingsunderlagSlide i React, route: /underlag)
\begin{frame}[label=bokforingsunderlag]
  \frametitle{Bokföringsunderlag och Bankinformation}
  \small
  För att analysera verksamhetskongruens och upptäcka privatkonsumtion behöver vi era fakturor och kontoutdrag.
  
  \vspace{0.3cm}
  \begin{block}{Varför behövs detta?}
    \footnotesize
    ✅ \textbf{Verksamhetskongruens:} Vi kontrollerar att inköp matchar verksamhetsbeskrivningen\\
    \hspace{0.5cm}\textit{Exempel: Blästringsföretag ska inte köpa motorcykeldäck}\\
    ✅ \textbf{Privatkonsumtion:} Vi flaggar ICA, H\&M, Willys bokförda på företaget\\
    ✅ \textbf{Konkurskontroll:} Vi varnar om betalning till avregistrerade företag\\
    ✅ \textbf{Leveransadress-validering:} Vi jämför mot företagets registrerade adress
  \end{block}
  
  \vspace{0.3cm}
  \begin{block}{💡 Tips: Du kan ladda upp blandat!}
    \footnotesize
    Det är helt okej att ladda upp alla dokument i en zon, blandat. Vår AI kommer automatiskt att sortera och kategorisera dokumenten åt dig.\\
    \textbf{OBS:} Detta tar större datakraft och längre bearbetningstid (ca 2-5 min extra), men det är ett tillåtet alternativ om du har mycket data.
  \end{block}
  
  \vspace{0.3cm}
  \textbf{Bankinformation (valfritt):}\\
  \vspace{0.2cm}
  \begin{minipage}{0.48\textwidth}
    \footnotesize
    \textbf{Bankgiro eller Plusgiro:}\\
    \fbox{\rule{4.5cm}{0.5cm}}\\
    \tiny\textcolor{gray}{Format: 123-4567 (bankgiro) eller 12345-6 (plusgiro)}\\
    \tiny\textcolor{gray}{IBAN anges på nästa slide (Bokföringsdata)}
  \end{minipage}
  
  \vspace{0.4cm}
  \textbf{Ladda upp bokföringsunderlag (5 kategorier):}\\
  \vspace{0.2cm}
  \footnotesize
  \begin{itemize}
    \item \textbf{Kontoutdrag:} Bankkontoutdrag (PDF)
    \item \textbf{Leverantörsfakturor:} Inköpsfakturor från era leverantörer
    \item \textbf{Kundfakturor:} Försäljningsfakturor till era kunder
    \item \textbf{Kvitton:} Kontantkvitton och småutlägg
    \item \textbf{Momsrapporter:} XML eller PDF från Skatteverket
    \item \textbf{Årsredovisningar (saknade år):} PDF - vi hämtar digitalt insända från 2021 via Bolagsverket API (endast juridiska personer)
  \end{itemize}
  
  \vspace{0.3cm}
  \textbf{OBS:} SIE-filer, IBAN-nummer och inkomstdeklarationer (K10/INK2) hämtas programvarumässigt via Skatteverket på nästa slide (Bokföringsdata).
  
  \vspace{0.3cm}
  \fbox{\begin{minipage}{10cm}
    \centering
    \vspace{0.3cm}
    �\\
    \vspace{0.2cm}
    \footnotesize\textbf{Dra och släpp alla filer här (PDF/PNG/JPG/XML)}\\
    \vspace{0.1cm}
    \tiny{Du kan ladda upp blandat - vår AI sorterar automatiskt (tar 2-5 min extra)}\\
    \vspace{0.1cm}
    \tiny{Eller använd separata zoner för varje kategori för snabbare bearbetning}\\
    \vspace{0.3cm}
  \end{minipage}}
  
  \vspace{0.3cm}
  \begin{flushright}
    \hyperlink{likviditet}{\beamerbutton{Nästa}}
  \end{flushright}
\end{frame}

% ═══════════════════════════════════════════════════════════════════════
% EKONOMISK RÅDGIVNING (Slides 11-14)
% Värdeskapande analys som ges FÖRE eventuell djupgranskning/avslag
% ═══════════════════════════════════════════════════════════════════════

% Slide 11: Likviditetsanalys
\begin{frame}[label=likviditet]
  \frametitle{Likviditetsanalys}
  \small
  Baserat på bankkontots transaktionshistorik kan vi ge en översikt av företagets likviditetsutveckling.
  
  \vspace{0.5cm}
  \begin{center}
    \textbf{Graf: Likviditet över tid (12 månader)}\\
    \vspace{0.3cm}
    \begin{tikzpicture}[scale=0.8]
      % Enkel linjegraf - ersätts med Recharts i React
      \draw[->] (0,0) -- (10,0) node[right] {Månad};
      \draw[->] (0,0) -- (0,4) node[above] {Belopp (KSEK)};
      \draw[thick, brand] plot coordinates {(0,2) (2,2.3) (4,1.8) (6,2.5) (8,2.8) (10,3.2)};
      \node at (5,-1) {\footnotesize Jan → Dec 2024};
    \end{tikzpicture}
  \end{center}
  
  \vspace{0.5cm}
  \begin{minipage}{0.48\textwidth}
    \textbf{Nuvarande likviditet:} 320 000 kr\\
    \textbf{Genomsnitt (12 mån):} 250 000 kr\\
    \textbf{Trend:} \textcolor{green}{Positiv (+28\%)}
  \end{minipage}
  \hfill
  \begin{minipage}{0.48\textwidth}
    \textbf{AI-analys:}\\
    \footnotesize
    ✅ Inga layering-mönster\\
    ✅ Normala transaktionsstorlekar\\
    ⚠️ 2 ovanligt stora insättningar (granskas)
  \end{minipage}
  
  \vspace{0.5cm}
  \begin{block}{Rekommendation}
    Byråchefen kan ge kommentarer kring likviditetsutvecklingen och föreslå åtgärder vid behov.
  \end{block}
  
  \vspace{0.5cm}
  \begin{flushright}
    \hyperlink{nextslide}{\beamerbutton{Nästa}}
  \end{flushright}
\end{frame}

% Slide 12: Omsättningsanalys
\begin{frame}[label=omsattning]
  \frametitle{Omsättningsanalys}
  \small
  Från bokföringens resultatrapporter kan vi se hur omsättningen utvecklats över åren.
  
  \vspace{0.5cm}
  \begin{center}
    \textbf{Graf: Omsättning per räkenskapsår}\\
    \vspace{0.3cm}
    \begin{tikzpicture}[scale=0.8]
      % Stapeldiagram - ersätts med Recharts BarChart i React
      \draw[->] (0,0) -- (10,0) node[right] {År};
      \draw[->] (0,0) -- (0,4) node[above] {Omsättning (MSEK)};
      \fill[brand!70] (0.5,0) rectangle (1.5,1.5);
      \fill[brand!70] (2.5,0) rectangle (3.5,2.0);
      \fill[brand!70] (4.5,0) rectangle (5.5,2.4);
      \fill[brand!70] (6.5,0) rectangle (7.5,2.8);
      \fill[brand!70] (8.5,0) rectangle (9.5,3.2);
      \node at (1,-0.5) {\footnotesize 2020};
      \node at (3,-0.5) {\footnotesize 2021};
      \node at (5,-0.5) {\footnotesize 2022};
      \node at (7,-0.5) {\footnotesize 2023};
      \node at (9,-0.5) {\footnotesize 2024};
    \end{tikzpicture}
  \end{center}
  
  \vspace{0.5cm}
  \begin{minipage}{0.48\textwidth}
    \textbf{Omsättning 2024:} 3,2 MSEK\\
    \textbf{Tillväxt (YoY):} \textcolor{green}{+14\%}\\
    \textbf{CAGR (5 år):} +16\%
  \end{minipage}
  \hfill
  \begin{minipage}{0.48\textwidth}
    \textbf{Analys:}\\
    \footnotesize
    ✅ Stabil tillväxt\\
    ✅ Inga ovanliga hopp\\
    ℹ️ Liten nedgång 2021 (pandemi?)
  \end{minipage}
  
  \vspace{0.5cm}
  \begin{block}{Rekommendation}
    Omsättningstillväxten är sund och följer en naturlig kurva. Fortsätt fokusera på tillväxtstrategier.
  \end{block}
  
  \vspace{0.5cm}
  \begin{flushright}
    \hyperlink{nextslide}{\beamerbutton{Nästa}}
  \end{flushright}
\end{frame}

% Slide 13: Resultatanalys
\begin{frame}[label=resultat]
  \frametitle{Resultatanalys}
  \small
  Från balansrapporterna kan vi se lönsamhetsutvecklingen över tid.
  
  \vspace{0.5cm}
  \begin{center}
    \textbf{Graf: Årets resultat per räkenskapsår}\\
    \vspace{0.3cm}
    \begin{tikzpicture}[scale=0.8]
      % Linjediagram med områdesfyllning - ersätts med Recharts AreaChart
      \draw[->] (0,0) -- (10,0) node[right] {År};
      \draw[->] (0,0) -- (0,4) node[above] {Resultat (KSEK)};
      \fill[brand!30] (0,0) -- (0,1.2) -- (2,1.5) -- (4,1.0) -- (6,1.8) -- (8,2.2) -- (10,2.5) -- (10,0) -- cycle;
      \draw[thick, brand] plot coordinates {(0,1.2) (2,1.5) (4,1.0) (6,1.8) (8,2.2) (10,2.5)};
      \node at (5,-1) {\footnotesize 2020 → 2024};
    \end{tikzpicture}
  \end{center}
  
  \vspace{0.5cm}
  \begin{minipage}{0.48\textwidth}
    \textbf{Resultat 2024:} 250 000 kr\\
    \textbf{Rörelsemarginal:} 7,8\%\\
    \textbf{Trend:} \textcolor{green}{Förbättras}
  \end{minipage}
  \hfill
  \begin{minipage}{0.48\textwidth}
    \textbf{Analys:}\\
    \footnotesize
    ✅ Positiva resultat alla år\\
    ℹ️ Lägre 2022 (högre kostnader?)\\
    ✅ Stark återhämtning 2023-2024
  \end{minipage}
  
  \vspace{0.5cm}
  \begin{block}{Rekommendation}
    Företaget är lönsamt och visar god utveckling. Överväg att optimera kostnadsstrukturen ytterligare.
  \end{block}
  
  \vspace{0.5cm}
  \begin{flushright}
    \hyperlink{nextslide}{\beamerbutton{Nästa}}
  \end{flushright}
\end{frame}

% Slide 14: Branschjämförelse
\begin{frame}[label=bransch]
  \frametitle{Branschjämförelse mot SCB-data}
  \small
  Hur står sig företaget mot branschgenomsnittet?
  
  \vspace{0.5cm}
  \begin{center}
    \textbf{Nyckeltal: Företaget vs Bransch (SNI 62010)}\\
    \vspace{0.3cm}
    \begin{tabular}{l|c|c|c}
      \textbf{Nyckeltal} & \textbf{Företaget} & \textbf{Bransch (SCB)} & \textbf{Status} \\ \hline
      Rörelsemarginal & 7,8\% & 8,5\% & ⚠️ Något lägre \\
      Soliditet & 42\% & 35\% & ✅ Bättre \\
      Kassalikviditet & 180\% & 150\% & ✅ Bättre \\
      Omsättning/anställd & 1,6 MSEK & 1,4 MSEK & ✅ Bättre \\
      Skuldsättningsgrad & 1,4 & 1,9 & ✅ Lägre (bättre) \\
    \end{tabular}
  \end{center}
  
  \vspace{0.5cm}
  \begin{block}{Sammanfattning}
    \footnotesize
    \textbf{Styrkor:}\\
    • Stark soliditet och kassalikviditet ger god finansiell stabilitet\\
    • Låg skuldsättning minskar finansiell risk\\
    • Hög produktivitet (omsättning/anställd)
    
    \vspace{0.3cm}
    \textbf{Förbättringsområden:}\\
    • Rörelsemarginalen kan förbättras genom kostnadsoptimering\\
    • Analysera varför marginalen ligger något under branschsnittet
  \end{block}
  
  \vspace{0.5cm}
  \begin{flushright}
    \hyperlink{nextslide}{\beamerbutton{Nästa}}
  \end{flushright}
\end{frame}

% ═══════════════════════════════════════════════════════════════════════
% DJUPGRANSKNING (Slides 15-16)
% Detaljerad analys som kan leda till fördjupad kontroll eller avslag
% ═══════════════════════════════════════════════════════════════════════

% Slide 19: Bokföringsanalys och verifikationsfel (BokforingsanalysSlide i React, route: /bokanalys)
\begin{frame}[label=bokanalys]
  \frametitle{Bokföringsanalys och verifikationsfel}
  \small
  AI-baserad granskning av verifikationer från SIE-filer. Endast flaggade poster visas.
  
  \vspace{0.4cm}
  \begin{minipage}{0.48\textwidth}
    \textbf{Sammanfattning:}\\
    \vspace{0.2cm}
    \begin{itemize}
      \item \textbf{4} Flaggade poster
      \item \textbf{2} Allvarliga avvikelser
      \item \textbf{2} Varningar
      \item \textbf{25\,732} Godkända
    \end{itemize}
  \end{minipage}
  \hfill
  \begin{minipage}{0.48\textwidth}
    \textbf{AI-granskning:}\\
    \vspace{0.2cm}
    \footnotesize
    Samtliga verifikationer har analyserats automatiskt. Poster som avviker från normala mönster flaggas för manuell granskning.
  \end{minipage}
  
  \vspace{0.5cm}
  \begin{block}{Flaggade verifikationer (klickbara rader)}
    \footnotesize
    \begin{tabular}{|l|l|l|p{4.5cm}|}
      \hline
      \textbf{Ver.nr} & \textbf{Datum} & \textbf{Belopp} & \textbf{Beskrivning} \\
      \hline
      A1 & 2024-01-15 & 850\,000 kr & Ovanligt hög konsultkostnad \\
      \hline
      A2 & 2024-03-22 & 45\,000 kr & Kontantuttag utan kvitto \\
      \hline
      A5 & 2024-09-20 & 8\,500 kr & Reseersättning kontant \\
      \hline
      A6 & 2024-10-12 & 320\,000 kr & Konsulttjänster från Cypern \\
      \hline
    \end{tabular}
  \end{block}
  
  \vspace{0.3cm}
  \begin{block}{Funktionalitet}
    \footnotesize
    \textbf{Klick på rad:} Öppnar modal med fullständig verifikation (debet/kredit, belopp, AI-flaggorsak)\\
    \textbf{Bedömning:} Användaren avgör om avvikelsen är acceptabel eller om fördjupad kontroll krävs
  \end{block}
  
  \vspace{0.5cm}
  \begin{flushright}
    \hyperlink{fortnoxkoppling}{\beamerbutton{Nästa}}
  \end{flushright}
\end{frame}



% Slide 20: Riskbedömning och besked (RiskbedomningSlide i React, route: /beslut)
\begin{frame}[label=riskbesked]
  \frametitle{Riskbedömning och besked}
  \small
  Efter analys av insamlad data görs en samlad riskbedömning baserat på:
  
  \vspace{0.3cm}
  \begin{itemize}
    \footnotesize
    \item Företagsdata och ägarstruktur
    \item Ekonomisk analys (likviditet, omsättning, resultat, branschjämförelse)
    \item Bokföringsanalys (verifikationer, avvikelser)
    \item Riskindikatorer (sanktionslistor, PEP, högriskländer)
  \end{itemize}
  
  \vspace{0.4cm}
  \begin{block}{Beslut (Endast för byråchef)}
    \footnotesize
    \textbf{Välj besked:}\\
    \vspace{0.2cm}
    $\square$ \textbf{Acceptera kund} – Gå vidare till avtal\\
    $\square$ \textbf{Fördjupad kontroll} – Begär ytterligare dokumentation\\
    $\square$ \textbf{Avslag} – Kunden accepteras ej\\
  \end{block}
  
  \vspace{0.4cm}
  \begin{block}{Prissättning (vid Accept)}
    \footnotesize
    Baserat på kundens valda tjänster och behov:
    \vspace{0.3cm}
    
    \textbf{Månadspris:} \underline{\hspace{4cm}} SEK \textit{(exklusive moms)}\\
    \vspace{0.2cm}
    
    \textit{Detta pris kommer att framgå i avtalet som genereras på nästa slide.}\\
    \textit{Prisförslag baseras på: Löpande bokföring, deklarationshjälp, fakturering, rådgivning, etc.}
  \end{block}
  
  \vspace{0.4cm}
  \begin{block}{AI-rekommendation}
    \footnotesize
    AI-systemet har analyserat kundens risk och komplexitet.\\
    \textbf{Rekommenderat pris:} \underline{\hspace{3cm}} SEK/mån (exkl. moms)
  \end{block}
  
  \vspace{0.4cm}
  \begin{flushright}
    \hyperlink{skyldigheter}{\beamerbutton{Generera avtal \& Gå vidare}}
  \end{flushright}
\end{frame}

% Slide 21: Kundens förväntade skyldigheter (SkyldigheterSlide i React, route: /skyldigheter)
\begin{frame}[label=skyldigheter]
  \frametitle{Kundens förväntade skyldigheter}
  \small
  Efter att kunden godkänts i riskbedömningen informeras om vilka skyldigheter som gäller under samarbetet:
  \begin{itemize}
    \item Faktura och förenklad faktura - korrekt utseende enligt Skatteverkets regler.
    \item Arkivering av bokföringsmaterial i minst sju år.
    \item Lämna korrekta, fullständiga och aktuella uppgifter till byrån.
    \item Inkomma med begärda underlag i tid.
    \item Omgående rapportera förändringar i verksamheten eller annan relevant information.
    \item Samarbeta vid frågor om kundkännedom och verksamhetens art.
    \item Förstå och följa grundläggande redovisningsskyldigheter.
    \item Vara medveten om att samarbetet kan avslutas vid avtalsbrott eller bristande efterlevnad.
  \end{itemize}
  \vspace{0.5cm}
  \textit{Dessa skyldigheter regleras i uppdragsavtalet och enligt lag.}
  \vspace{0.8cm}
  \begin{flushright}
    \beamerbutton{Nästa}
  \end{flushright}
\end{frame}

% Slide 22: Avtalsvillkor och godkännande (AvtalSlide i React, route: /avtal)
% BankID-signering med PDF-visning och sidebar
\begin{frame}[fragile]
  \frametitle{Avtalsvillkor och godkännande}

  \begin{block}{Avtal (skrolla för att läsa)}
    \vspace{0.2cm}
    % Här visas avtalet i en inbäddad PDF-komponent.
    % Exempelvis:
    % \includepdf[pages=1,width=\linewidth,pagecommand={}]{avtal.pdf}
    %
    % I webbgränssnittet motsvarar detta en skrollbar vy.
    \textit{\footnotesize Avtalet visas här i PDF-format. Användaren kan skrolla för att läsa hela innehållet.}
    \vspace{0.4cm}
  \end{block}

  \vspace{0.4cm}
  \begin{center}
    \beamerbutton{Signera med BankID}
  \end{center}

  \vspace{0.6cm}
  \small
  Efter genomförd signering ersätts PDF-dokumentet automatiskt med en version där
  namnet från BankID-identifieringen infogas som signatur i dokumentet.
  \medskip

  \textit{Exempel:}\\
  \textbf{Signerad av:} Anna Svensson (BankID)

\end{frame}

% Slide 23: E-postbekräftelse och leverans av dokument (DocumentDeliverySlide i React, route: /dokument)
\begin{frame}[fragile]
  \frametitle{Bekräftelse och leverans av material}

  \begin{block}{Ange e-postadress för kopia av dokument}
    \vspace{0.2cm}
    \textit{För din dokumentation skickas en kopia av alla relevanta filer till den e-postadress du anger nedan.}
    \vspace{0.4cm}

    % Fältet motsvarar en e-postinput i appens gränssnitt
    \textbf{E-postadress:} \underline{\hspace{6cm}}
    \vspace{0.5cm}

    \begin{itemize}
      \item Avtal (PDF med BankID-signatur)
      \item Riskanalys och grafisk rapport (PDF)
      \item Sammanställning av dina svar
    \end{itemize}
  \end{block}

  \vspace{0.5cm}
  \begin{center}
    \beamerbutton{Skicka material till min e-post}
  \end{center}

  \vspace{0.6cm}
  \small
  \textit{När du bekräftar skickas all information krypterat via säker anslutning.  
  Du får även ett meddelande som bekräftar att signeringen sVälkommen som kund! 🎉lutförts.}

\end{frame}

%%%%%%%%%%%%%%%%%%%%%%%%%%%%%%%%%%%%%%%%%%%%%%%%%%%
% POST-KONTRAKT SETUP: Hjälp kunden komma igång
%%%%%%%%%%%%%%%%%%%%%%%%%%%%%%%%%%%%%%%%%%%%%%%%%%%

% Slide 23: Koppling till Fortnox – Välja rätt paket (FortnoxPackageSlide i React, route: /fortnox)
\begin{frame}[label=fortnoxkoppling]
  \frametitle{Koppling till Fortnox – Ditt bokföringsprogram}
  \small
  Vi använder Fortnox för löpande bokföring. Du behöver ett eget Fortnox-abonnemang.\\
  \textit{"Vi rekommenderar våra klienter att köpa Mini eller Mellan. Om du gör beställningen med en gång kan jag koppla upp dig direkt till vårt företagskonto."}
  
  \vspace{0.3cm}
  \begin{block}{Rekommenderade paket (För enmansföretagare)}
    \footnotesize
    \begin{tabular}{|l|l|p{5.5cm}|r|}
      \hline
      \textbf{Paket} & \textbf{För vem?} & \textbf{Inkluderar} & \textbf{Pris} \\
      \hline
      Mini & Enklast & Bokföring, bank, fakturering, företagskort & 189 kr/mån \\
      \hline
      \textbf{Liten} & \textbf{Populärast} & \textbf{+ Lön till dig själv, integrationer} & \textbf{319 kr/mån} \\
      \hline
      \textbf{Mellan} & \textbf{Rekommenderat} & \textbf{+ Avancerade rapporter} & \textbf{489 kr/mån} \\
      \hline
      Stor & Mest funktioner & + Budgetverktyg, analysverktyg & 709 kr/mån \\
      \hline
    \end{tabular}
  \end{block}
  
  \vspace{0.2cm}
  \begin{block}{För dig med anställda (+ paketen)}
    \footnotesize
    \textbf{Mini+} (329 kr/mån), \textbf{Liten+} (439 kr/mån), \textbf{Mellan+} (589 kr/mån), \textbf{Stor+} (839 kr/mån)\\
    Inkluderar: Lönehantering, faktura-attestering, anläggningsregister
  \end{block}
  
  \vspace{0.2cm}
  \begin{block}{Erbjudande för nystartade företag}
    \footnotesize
    \textbf{Kod: NYSTARTAD} – Valfritt paket kostnadsfritt i 6 månader!\\
    \textit{(Gäller företag startade inom senaste 3 månaderna)}
  \end{block}
  
  \vspace{0.3cm}
  \begin{center}
    \hyperlink{https://www.fortnox.se/paket}{\beamerbutton{Beställ Fortnox-paket}}
  \end{center}
  
  \vspace{0.3cm}
  \begin{flushright}
    \hyperlink{bankrattigheter}{\beamerbutton{Nästa: Bankkoppling}}
  \end{flushright}
\end{frame}

% Slide 24: Läsrättigheter till företagsbankkonto (BankRattigheterSlide i React, route: /bank)
\begin{frame}[fragile,label=bankrattigheter]
  \frametitle{Steg 1: Läsrättigheter till företagets bankkonto}

  \begin{block}{Varför behövs detta?}
    \vspace{0.2cm}
    För att vi ska kunna hjälpa dig med löpande bokföring behöver vi läsrättigheter till ditt företagsbankkonto.
    Detta gör att vi automatiskt kan hämta kontoutdrag och matcha betalningar mot fakturor.
    \vspace{0.2cm}
    \textbf{OBS:} Vi får ENDAST läsrättigheter – vi kan aldrig göra betalningar eller överföringar.
  \end{block}

  \vspace{0.4cm}

  \begin{block}{Så här gör du (välj din bank):}
    \footnotesize
    \begin{itemize}
      \item \textbf{SEB:} Logga in på företag.seb.se → Inställningar → Behörigheter → Lägg till läsrättighet → Ange vårt org.nr: [AUTO-IFYLLT]
      \item \textbf{Handelsbanken:} Företag.handelsbanken.se → Behörigheter → Kontoåtkomst → Lägg till extern part → Org.nr: [AUTO-IFYLLT]
      \item \textbf{Nordea:} Logga in på företag.nordea.se → Administrera → Behörigheter → Open Banking → Ge läsrättighet → [AUTO-IFYLLT]
      \item \textbf{Swedbank:} Företag.swedbank.se → Inställningar → Företagsbehörigheter → Lägg till mottagare → [AUTO-IFYLLT]
      \item \textbf{Andra banker:} Kontakta din bank och be om läsrättighet (PSD2/Open Banking) för org.nr: [AUTO-IFYLLT]
    \end{itemize}
  \end{block}

  \vspace{0.4cm}
  \begin{center}
    \beamerbutton{✓ Klart – jag har gett läsrättighet}
  \end{center}

  \vspace{0.3cm}
  \footnotesize
  \textit{Tips: Ta en skärmdump när du givit behörigheten – då har du ett kvitto för framtida referens.}

\end{frame}

% ----------------------------------------------------------
% Slide 25: Deklarationsombud via Skatteverket (DeklarationsombudSlide i React, route: /ombud)
\begin{frame}[fragile]
  \frametitle{Steg 2: Lägg till oss som deklarationsombud}

  \begin{block}{Varför behövs detta?}
    \vspace{0.2cm}
    För att vi ska kunna deklarera F-skatt, moms och arbetsgivaravgifter åt dig behöver du lägga till oss
    som deklarationsombud hos Skatteverket.
  \end{block}

  \vspace{0.4cm}

  \begin{block}{Så här gör du:}
    \footnotesize
    \begin{enumerate}
      \item Gå till \textbf{skatteverket.se/etjanster}
      \item Logga in med BankID
      \item Välj \textbf{"Mina sidor"} → \textbf{"Fullmakter"} → \textbf{"Lägg till ombud"}
      \item Ange vårt organisationsnummer: \textbf{[AUTO-IFYLLT från config.json]}
      \item Välj behörigheter:
        \begin{itemize}
          \item ✓ Deklarationsombud (F-skatt)
          \item ✓ Momsdeklaration
          \item ✓ Arbetsgivardeklaration
          \item ✓ Inkomstdeklaration för företag
        \end{itemize}
      \item Bekräfta med BankID
    \end{enumerate}
  \end{block}

  \vspace{0.4cm}
  \begin{center}
    \beamerbutton{Verifiera att jag lagt till ombud}
  \end{center}

  \vspace{0.3cm}
  \footnotesize
  \textit{När du klickar verifierar vi via Skatteverkets Ombudshantering API att behörigheten är registrerad.}\\
  \textit{Om det inte syns än: Det kan ta några minuter – prova igen om en stund.}

\end{frame}

% ----------------------------------------------------------
% Slide 26: Digital dokumenthantering (Scan to Dropbox/Email) - EJ IMPLEMENTERAD ÄN
% TODO: Skapa DocumentSetupSlide.jsx för route: /dokument-setup
\begin{frame}[fragile]
  \frametitle{Steg 3: Sätt upp digital dokumenthantering}

  \begin{block}{Varför är detta viktigt?}
    \vspace{0.2cm}
    För att vi ska kunna bokföra snabbt och korrekt behöver vi få dina underlag digitalt.
    Med rätt setup behöver du bara lägga fakturor och kvitton i skannern – resten sker automatiskt!
  \end{block}

  \vspace{0.4cm}Bokförda transaktioner

Sök bland transaktioner som är bokförda till och med 2025-10-20 kl. 03.03

Om du väljer att visa alla typer av transaktioner kan du även beställa utskrift av sökresultatet. Utskriften skickar vi till dig eller till den mottagare som du väljer.
Val av tidsperiod Val av tidsperiod
Valfri period
Datum från och med Datum från och med
Skriv datum i formatet år fyra siffror, månad två siffror och dag två siffror eller välj ur kalendern.

Måndag 1 januari 2018
Datum till och med Datum till och med
Skriv datum i formatet år fyra siffror, månad två siffror och dag två siffror eller välj ur kalendern.

Måndag 20 oktober 2025
Typ av transaktion Typ av transaktion
Välj vilka typer av kontohändelser eller transaktioner som du vill söka fram.

Välj "Alla" om du vill beställa utskrift av sökresultatet (kontohistorik).
transaktionstypDropdownBeskrivningtransaktionstypDropdownBeskrivningExtra
Välj typ (Alla valda)
Sökresultat (201)


  \begin{block}{Alternativ 1: Scan to Email (enklast)}
    \footnotesize
    \textbf{Steg:}
    \begin{enumerate}
      \item Hitta "Scan to Email"-funktionen i din skanner (oftast under Settings/Network)
      \item Lägg till vår e-postadress: \textbf{underlag@[DIN-BYRA].se} [AUTO-IFYLLT]
      \item Testa genom att scanna ett kvitto och skicka
      \item Vi bekräftar när vi mottagit test-dokumentet ✓
    \end{enumerate}Välkommen som kund! 🎉
    \textbf{Funkar med:} De flesta kontorsskannrar (Fujitsu, Brother, HP, Canon, Epson)
  \end{block}

  \vspace{0.4cm}

  \begin{block}{Alternativ 2: Scan to Dropbox (rekommenderas)}
    \footnotesize
    \textbf{Steg:}
    \begin{enumerate}
      \item Vi skapar en delad Dropbox-mapp: \textbf{[FÖRETAGSNAMN]\_Underlag}
      \item Du får en inbjudan på e-post – acceptera den
      \item Konfigurera din skanner att spara till Dropbox-mappen
      \item Skapa undermappar: \texttt{Leverantörsfakturor}, \texttt{Kvitton}, \texttt{Avtal}, \texttt{Löner}
    \end{enumerate}
    \textbf{Fördel:} Du kan släppa filer från mobilen, datorn eller skannern – allt hamnar på samma ställe!
  \end{block}

  \vspace{0.4cm}
  \begin{center}
    \beamerbutton{✓ Klart – jag har konfigurerat dokumenthantering}
  \end{center}

\end{frame}

% ----------------------------------------------------------
% SLIDE 24b: Guide för vanliga skannermärken (expanderbar)
% ----------------------------------------------------------
\begin{frame}[fragile]
  \frametitle{Steg 3b: Scannerguide (per märke)}
  \framesubtitle{Expandera din scannermodell för detaljerad guide}

  \begin{block}{Fujitsu ScanSnap (iX-serien, S1300, etc.)}
    \footnotesize
    \begin{enumerate}
      \item Öppna \textbf{ScanSnap Manager}
      \item Högerklicka på ikonen → \textbf{Settings}
      \item Under "Save" välj \textbf{"Specify a folder"}
      \item Bläddra till din Dropbox-mapp: \texttt{Dropbox/[FÖRETAGSNAMN]\_Underlag/}
      \item Klart! Testa genom att scanna ett dokument
    \end{enumerate}
  \end{block}

  \vspace{0.3cm}

  \begin{block}{Brother (ADS-serien, MFC-serien)}
    \footnotesize
    \begin{enumerate}
      \item Tryck på \textbf{Scan} på skannerns display
      \item Välj \textbf{"to E-mail"} eller \textbf{"to Network"}
      \item Lägg till e-post: \textbf{underlag@[DIN-BYRA].se}
      \item Spara inställningen som favorit (t.ex. "Bokföring")
      \item Framöver: Tryck bara på favoriten och scanna!
    \end{enumerate}
  \end{block}

  \vspace{0.3cm}

  \begin{block}{HP OfficeJet / LaserJet (med ADF)}
    \footnotesize
    \begin{enumerate}
      \item Gå till skannerns webbgränssnitt (skriv IP-adressen i webbläsare)
      \item Välj \textbf{Scan/Digital Send} → \textbf{Email}
      \item Konfigurera SMTP (vi hjälper dig med detta)
      \item Lägg till mottagare: \textbf{underlag@[DIN-BYRA].se}
      \item Spara som snabbval
    \end{enumerate}
  \end{block}

  \vspace{0.3cm}
  \footnotesize
  \textit{Behöver du hjälp? Ring oss på [TELEFONNUMMER] så går vi igenom setup tillsammans! (5-10 min)}

\end{frame}


% Slide 27: Välkommen som kund! 🎉 [WelcomeSlide.jsx]
% React implementation: /src/components/Slides/WelcomeSlide.jsx
% Route: /welcome - Final onboarding summary with completed steps, next steps, and contact info
\begin{frame}[fragile]
  \frametitle{🎉 Välkommen som kund!}

  \begin{block}{Grattis – du är nu redo att börja!}
    \vspace{0.2cm}
    Alla steg är nu klara och vi kan börja arbeta med din bokföring.
    Här är en sammanfattning av vad som hänt:
  \end{block}

  \vspace{0.4cm}

  \begin{block}{✓ Klart:}
    \footnotesize
    \begin{itemize}
      \item ✓ Du har genomgått vår onboarding-process
      \item ✓ Riskbedömning genomförd och godkänd
      \item ✓ Avtal signerat med BankID
      \item ✓ Läsrättigheter till bankkonto uppsatta
      \item ✓ Deklarationsombud registrerat hos Skatteverket
      \item ✓ Digital dokumenthantering konfigurerad
    \end{itemize}
  \end{block}

  \vspace{0.4cm}

  \begin{block}{Nästa steg:}
    \footnotesize
    \begin{enumerate}
      \item \textbf{Inom 24 timmar:} Din personliga redovisningskonsult kontaktar dig för att boka uppstartsmöte
      \item \textbf{Uppstartsmöte:} Vi går igenom din verksamhet, förväntningar och bokföringsrutiner (30-45 min)
      \item \textbf{Första bokföringen:} Vi påbörjar arbetet med din löpande bokföring
      \item \textbf{Månadsmöte:} Varje månad får du en ekonomisk rapport och rådgivning
    \end{enumerate}
  \end{block}

  \vspace{0.4cm}



  \vspace{0.4cm}
  \begin{center}
    \Large\textbf{Tack för ditt förtroende!}\\
    \vspace{0.2cm}
    \normalsize\textit{Vi ser fram emot ett långt och framgångsrikt samarbete.}
  \end{center}

  \vspace{0.5cm}
  \begin{flushright}
    \beamerbutton{Nästa: Löpande rutiner}
  \end{flushright}

\end{frame}

%%%%%%%%%%%%%%%%%%%%%%%%%%%%%%%%%%%%%%%%%%%%%%%%%%%

% Slide 28: Löpande rutiner [OngoingRoutinesSlide.jsx]
% React implementation: /src/components/Slides/OngoingRoutinesSlide.jsx
% Route: /rutiner - KYC updates, risk reassessment, document submission schedules
\begin{frame}[fragile]
  \frametitle{Löpande rutiner}
  \framesubtitle{Så här håller vi din verksamhet uppdaterad och regelkonform}

  \begin{block}{Viktigt att veta}
    Du behöver inte komma ihåg alla dessa datum – vi skickar automatiska påminnelser och 
    sköter de flesta rutinerna åt dig. Din enda uppgift är att leverera underlag enligt 
    överenskommen tidsplan.
  \end{block}

  \vspace{0.3cm}

  \begin{block}{🔄 KYC-uppdateringar (Årligen)}
    \footnotesize
    \begin{itemize}
      \item Bekräfta ägarstruktur (förändringar i ägande över 25\%)
      \item Uppdatera kontaktinformation för verkliga huvudmän
      \item Verifiera verksamhetsbeskrivning och SNI-kod
      \item Granska eventuella förändringar i styrelsens sammansättning
    \end{itemize}
    \textit{Vi skickar ut formulär via e-post 30 dagar innan årsdagen}
  \end{block}

  \vspace{0.3cm}

  \begin{block}{🛡️ Årlig riskbedömning (Årligen)}
    \footnotesize
    \begin{itemize}
      \item Genomgång av transaktionsmönster från det senaste året
      \item Kontroll mot uppdaterade sanktionslistor (EU, UN, OFAC)
      \item Bedömning av PEP-status (Politically Exposed Persons)
      \item Dokumentation av riskklassificering (låg/medel/hög)
    \end{itemize}
    \textit{Automatisk körning varje år + vid väsentliga förändringar}
  \end{block}

  \vspace{0.3cm}

  \begin{block}{📄 Dokumentinlämning (Månatligen)}
    \footnotesize
    \begin{itemize}
      \item Senast den 10:e: Leverantörsfakturor för föregående månad
      \item Senast den 10:e: Kvitton och utlägg
      \item Vid löneperiod: Löneunderlag senast 5 dagar före utbetalning
    \end{itemize}
    \textit{Automatiska påminnelser 3 dagar innan deadline}
  \end{block}Bokförda transaktioner

Sök bland transaktioner som är bokförda till och med 2025-10-20 kl. 03.03

Om du väljer att visa alla typer av transaktioner kan du även beställa utskrift av sökresultatet. Utskriften skickar vi till dig eller till den mottagare som du väljer.
Val av tidsperiod Val av tidsperiod
Valfri period
Datum från och med Datum från och med
Skriv datum i formatet år fyra siffror, månad två siffror och dag två siffror eller välj ur kalendern.

Måndag 1 januari 2018
Datum till och med Datum till och med
Skriv datum i formatet år fyra siffror, månad två siffror och dag två siffror eller välj ur kalendern.

Måndag 20 oktober 2025
Typ av transaktion Typ av transaktion
Välj vilka typer av kontohändelser eller transaktioner som du vill söka fram.

Välj "Alla" om du vill beställa utskrift av sökresultatet (kontohistorik).
transaktionstypDropdownBeskrivningtransaktionstypDropdownBeskrivningExtra
Välj typ (Alla valda)
Sökresultat (201)


  \vspace{0.3cm}

  \begin{block}{📅 Compliance-kalender (Enligt schema)}
    \footnotesize
    \begin{itemize}
      \item Månatligen: Momsredovisning (vid månadsredovisning)
      \item Kvartalsvis: Arbetsgivardeklaration
      \item Maj varje år: Inkomstdeklaration och årsredovisning
    \end{itemize}
    \textit{Vi sköter alla inlämningar automatiskt – du får kopia}
  \end{block}

  \vspace{0.5cm}
  \begin{flushright}
    \beamerbutton{Nästa: Support \& kontakt}
  \end{flushright}

\end{frame}

%%%%%%%%%%%%%%%%%%%%%%%%%%%%%%%%%%%%%%%%%%%%%%%%%%%

% Slide 29: Support & Kontakt [SupportSlide.jsx]
% React implementation: /src/components/Slides/SupportSlide.jsx
% Route: /support - How to get help, escalation, compliance questions
\begin{frame}[fragile]Bokförda transaktioner

Sök bland transaktioner som är bokförda till och med 2025-10-20 kl. 03.03

Om du väljer att visa alla typer av transaktioner kan du även beställa utskrift av sökresultatet. Utskriften skickar vi till dig eller till den mottagare som du väljer.
Val av tidsperiod Val av tidsperiod
Valfri period
Datum från och med Datum från och med
Skriv datum i formatet år fyra siffror, månad två siffror och dag två siffror eller välj ur kalendern.

Måndag 1 januari 2018
Datum till och med Datum till och med
Skriv datum i formatet år fyra siffror, månad två siffror och dag två siffror eller välj ur kalendern.

Måndag 20 oktober 2025
Typ av transaktion Typ av transaktion
Välj vilka typer av kontohändelser eller transaktioner som du vill söka fram.

Välj "Alla" om du vill beställa utskrift av sökresultatet (kontohistorik).
transaktionstypDropdownBeskrivningtransaktionstypDropdownBeskrivningExtra
Välj typ (Alla valda)
Sökresultat (201)

  \frametitle{Support \& Kontakt}
  \framesubtitle{Vi finns här för att hjälpa dig – välj rätt kanal för snabbast svar}

  \begin{columns}[T]
  \column{0.48\textwidth}
  
  \begin{block}{💻 Teknisk support}
    \footnotesize
    Problem med inloggning, systemet eller filuppladdning\\
    \textbf{E-post:} support@dinbyra.se\\
    \textbf{Telefon:} 08-123 45 67\\
    \textbf{Svar:} Inom 4 timmar\\
    \textit{Mån-Fre 08:00-17:00}
  \end{block}

  \vspace{0.3cm}

  \begin{block}{📊 Bokföringsfrågor}
    \footnotesize
    Frågor om bokföring, verifikationer och rapporter\\
    \textbf{Kontakt:} Din tilldelade redovisningskonsult\\
    \textbf{Svar:} Inom 24 timmar\\
    \textit{Se kontaktinformation i välkomstbrevet}
  \end{block}

  \vspace{0.3cm}

  \begin{block}{⚖️ Regelverksfrågor}
    \footnotesize
    Penningtvätt, KYC, sanktioner\\
    \textbf{E-post:} compliance@dinbyra.se\\
    \textbf{Telefon:} 08-123 45 68\\
    \textbf{Svar:} Inom 24h (akut: 2h)\\
    \textit{Mån-Fre 09:00-16:00}
  \end{block}

  \column{0.48\textwidth}
  
  \begin{block}{🏛️ Skattefrågor}
    \footnotesize
    Deklarationer, moms, Skatteverket\\
    \textbf{E-post:} skatt@dinbyra.se\\
    \textbf{Telefon:} 08-123 45 69\\
    \textbf{Svar:} Inom 24h (deadline: samma dag)\\
    \textit{Mån-Fre 08:00-17:00}
  \end{block}

  \vspace{0.3cm}

  \begin{block}{🚨 Akuta ärenden}
    \footnotesize
    Brådskande situationer\\
    \textbf{Ring alltid:} 08-123 45 70\\
    \textbf{E-post:} akut@dinbyra.se\\
    \textbf{Svar:} Omedelbar hantering\\
    \textit{Helgfria vardagar 08:00-17:00}
    
    \vspace{0.2cm}
    Exempel:
    \begin{itemize}
      \item Skatteverket kontaktar företaget
      \item Misstänkta bedrägerier
      \item Förelägganden med kort svarstid
    \end{itemize}
  \end{block}

  \end{columns}

  \vspace{0.4cm}

  \begin{block}{Eskaleringsprocess}
    \footnotesize
    \textbf{1.} Kontakta rätt supportkanal $\rightarrow$ 
    \textbf{2.} Får inget svar? Påminnelse efter 24h $\rightarrow$ 
    \textbf{3.} Akut? Ring 08-123 45 70
  \end{block}

  \vspace{0.3cm}

  \begin{block}{Självhjälpsresurser}
    \footnotesize
    \begin{itemize}
      \item \textbf{FAQ:} Vanliga frågor och snabba svar
      \item \textbf{Användarguider:} Steg-för-steg-instruktioner
      \item \textbf{Videotutorials:} Lär dig genom att se
      \item \textbf{Regelverksdokumentation:} AML, GDPR, bokföringslag
    \end{itemize}
  \end{block}

  \vspace{0.5cm}
  \begin{flushright}
    \beamerbutton{Klar med onboarding! 🎉}
  \end{flushright}

\end{frame}

%%%%%%%%%%%%%%%%%%%%%%%%%%%%%%%%%%%%%%%%%%%%%%%%%%%

%Sida för inställningar som erhålles vid klick på kuggghjulet uppe till höger

% --- Slide: Inställningar (via kugghjulet) ---
\begin{frame}[fragile]
  \frametitle{Inställningar – åtkomst via kugghjulet uppe till höger}

  \begin{block}{Syfte}
    Denna sida låter användaren hantera sitt konto och firmans grundinformation.
    För denna version används inga användarroller; allt hanteras av kontoägaren.
  \end{block}

  \vspace{0.4cm}

  \begin{block}{Funktionalitet}
    \begin{enumerate}
      \item \textbf{Firmakonfiguration} – Kontaktuppgifter som auto-fylls i dokument och avtal:
        \begin{itemize}
          \item Firmanamn, organisationsnummer
          \item Supportkontakt (e-post, telefon)
          \item Ansvarig för efterlevnad (compliance officer)
          \item Skattekontor och revisorsuppgifter
        \end{itemize}
      Dessa uppgifter används för att automatiskt fylla i kontaktinformation i avtal, rapporter och kommunikation med klienter.
      \vspace{0.2cm}
      
      \item \textbf{Ladda upp konfigurationsfil (config.json)} – \textit{Framtida funktion:}
        \begin{itemize}
          \item Textinnehåll för editerbara slides
          \item Antal år för SIE-filhämtning (standard: 7 år enligt lag, kan reduceras för nystartade företag)
          \item Valfria fält och textanpassningar
        \end{itemize}
      Vissa slides (t.ex.\ drop zones och inmatning av personnummer/org.nr) är låsta men kan ordnas om i listan.
      \textit{OBS:} För närvarande används config.json endast för avtalsdata.
      \vspace{0.2cm}
      
      \item \textbf{Användarhantering} – Hantera åtkomst till systemet:
        \begin{itemize}
          \item Registrera nya användare genom att skicka inbjudan till e-postadress
          \item Se lista över aktiva användare med roller och senaste inloggning
          \item \textbf{"Glömt lösenord"-funktionalitet} via e-post med säker återställningslänk
          \item Ta bort användare (kräver bekräftelse)
        \end{itemize}
      \vspace{0.2cm}
      
      \item \textbf{Risktester och validering} – Information om automatiska kontroller:
        \begin{itemize}
          \item Länk till fullständig dokumentation: \texttt{VALIDATION\_TESTS\_NEW.md}
          \item Översikt över testade områden:
            \begin{itemize}
              \item \textbf{Kategori 1:} SIE-filintegritet (KSUMMA, balansräkning, debet/kredit)
              \item \textbf{Kategori 2:} Bristande bokföringskompetens (okunniga bokförare)
              \item \textbf{Kategori 3:} Fuskindikationer (Test 3.3 Momsrapporter, 3.4 Avskrivningar, \textbf{3.7 Sammanslagna verifikationer, 3.8 Raderade transaktioner, 3.9 Sekvensbrott i verifikationsnummer})
              \item \textbf{Kategori 4:} Penningtvätt \& kundbedrägeri (layering, PEP, sanktionslistor)
            \end{itemize}
          \item \textbf{OBS:} Testerna kan \textit{inte} stängas av av användaren.
          \item Systemet "tjuter och plingar" vid avvikelser – detta är enligt lag obligatoriskt.
          \item Användare måste acceptera att alla valideringskontroller körs på all bokföringsdata.
          \item \textbf{Exempel på AI-detektering:} Test 3.7 använder Claude/GPT för att identifiera sammanslagna verifikationer (obfuscering) genom semantisk textanalys.
          \item \textbf{Test 3.8 - Raderade transaktioner:} Validerar att \texttt{\#BTRANS} (borttagna) och \texttt{\#RTRANS} (korrigerade) finns i SIE-filen enligt bokföringslagen. Flaggar export "Utan strukna rader" (Visma Administration).
          \item \textbf{Test 3.9 - Sekvensbrott:} Kontrollerar att verifikationsnummer är löpande inom varje serie (A1, A2, A3, ...). Saknade nummer indikerar manuell radering ur SIE-fil.
        \end{itemize}
      \vspace{0.2cm}
      
      \item \textbf{Prenumeration och fakturahantering:}
        \begin{itemize}
          \item Se aktuell prenumeration, status och förfallodatum
          \item Förnya prenumeration (dirigeras till betalning via Stripe/Klarna)
          \item Avsluta prenumeration (kräver bekräftelse, data raderas efter 30 dagar)
          \item Fakturahistorik via extern integration (Fortnox/Stripe)
          \item \textit{OBS:} Faktura-PDFer lagras INTE lokalt – endast metadata och externa länkar
          \item Kunden kan ladda ner fakturor direkt från faktureringstjänstens portal
        \end{itemize}
      \vspace{0.2cm}
      
      \item \textbf{Ta bort konto} – \textit{Danger Zone:}
        \begin{itemize}
          \item Permanent radering av alla kunddata och bokföringsunderlag
          \item Kräver dubbel bekräftelse (checkbox + bekräftelse-modal)
          \item Varning: All data raderas omedelbart och kan ej återställas
          \item Endast tillgängligt för kontoägare
        \end{itemize}
    \end{enumerate}
  \end{block}

  \vspace{0.4cm}

  \begin{block}{Framtida UI-idéer}
    \begin{itemize}
      \item React-komponenter i Tailwind CSS (brand green palette: brand-50 till brand-900)
      \item Moderna "card"-sektioner för varje inställningskategori med ikoner
      \item Tydlig visuell hierarki med \texttt{hover}-effekter och "expand/collapse"-komponenter
      \item Formulärvalidering i realtid för firmakonfiguration
      \item Integration med backend via REST API för användar- och filhantering
      \item Toast-notifikationer vid lyckade ändringar
    \end{itemize}
  \end{block}
\end{frame}
%%%%%%%%%%%%%%%%%%%%%%%%%%%%%%%%%%%%%%%%%%%%%%%%%%%


\begin{frame}{Administrationspanel - AdminDashboard}
\framesubtitle{Systemöversikt och hantering för Celestial Consulting AB}

\begin{block}{Syfte}
Ge systemadministratören (Lasse) realtidsöversikt och kontroll över plattformen utan att lagra personlig kunddata.
\end{block}

\begin{block}{1. Dashboard-översikt (Startsida)}
  \begin{itemize}
    \item \textbf{Statistikkort (Cards):}
      \begin{itemize}
        \item API-anrop: Idag / Denna månad (med trendpilar: \textcolor{green}{↗} stigande, \textcolor{red}{↘} fallande)
        \item Aktiva sessioner: Antal inloggade användare just nu (klickbar för att expandera lista)
        \item Totalt kunder: Antal registrerade kunder med tillväxt \% sedan förra månaden
        \item Systemdrift: Uptime \% (färgkodad: grön 99+\%, gul 95-99\%, röd <95\%)
      \end{itemize}
    \item \textbf{Recent Activity (Timeline):}
      \begin{itemize}
        \item Senaste 20 aktiviteter: Kundinloggning, dokumentuppladdning, risktest flaggad, supportärende skapat
        \item Tidsstämpel, ikon per aktivitetstyp, anonym kundidentifiering (customer\_hash)
        \item Klickbar för detaljer (öppnar modal med fullständig metadata)
      \end{itemize}
  \end{itemize}
\end{block}

\end{frame}

%%%%%%%%%%%%%%%%%%%%%%%%%%%%%%%%%%%%%%%%%%%%%%%%%%%

\begin{frame}{Administrationspanel - Support \& Fakturering}
\framesubtitle{Hantering av kundärenden och ekonomi}

\begin{block}{2. Support Queue (Ärendekö)}
  \begin{itemize}
    \item \textbf{Tabell med kolumner:}
      \begin{itemize}
        \item Kund: Anonym hash eller firmanamn (om tillåtet av kund)
        \item Ärendetyp: Dropdown (Teknisk, Bokföring, Compliance, Skatt, Brådskande)
        \item Prioritet: Badges med färgkodning (Hög=röd, Medel=orange, Låg=grön)
        \item Status: Open / In Progress / Resolved med färgkodning
        \item Väntetid: Beräknas automatiskt från skapelsetid (flagga om >24h)
        \item Åtgärder: Ikonknappar (Visa detaljer, Markera löst, Eskalera)
      \end{itemize}
    \item \textbf{Filter- och sorteringskontroller:}
      \begin{itemize}
        \item Filtrera på: Ärendetyp, Prioritet, Status
        \item Sortera efter: Väntetid (längst först), Prioritet, Skapelsedatum
      \end{itemize}
    \item \textbf{Ärendedetaljer (Modal):}
      \begin{itemize}
        \item Kundens meddelande, bifogade filer (screenshots), historik
        \item Textfält för svar, "Skicka svar"-knapp (triggar email via SendGrid)
        \item Intern anteckningsyta (syns ej för kund)
      \end{itemize}
  \end{itemize}
\end{block}

\begin{block}{3. Fakturering (Invoice Management)}
  \begin{itemize}
    \item \textbf{Fakturaöversikt (Tabell):}
      \begin{itemize}
        \item Kund, Belopp (kr), Status (Betald=grön, Väntande=orange, Förfallen=röd)
        \item Förfallodatum (highlight om överskridet), Fakturatyp (Prenumeration, Engångskostnad)
        \item Åtgärder: Ladda ner PDF (länk till Fortnox), Skicka påminnelse (email)
      \end{itemize}
    \item \textbf{Integration med Fortnox API:}
      \begin{itemize}
        \item Metadata lagras lokalt (invoice\_id, belopp, datum, status, fortnox\_url)
        \item PDF:er hämtas on-demand från Fortnox via proxy-route med temporär signerad URL
        \item "Skicka påminnelse" genererar automatiskt email med betalningslänk
      \end{itemize}
    \item \textbf{Mock-data för LIA-demo:}
      \begin{itemize}
        \item 10 testfakturor med varierande status (5 betalda, 3 väntande, 2 förfallna)
        \item Realistiska belopp (1 200 - 5 000 kr/månad)
      \end{itemize}
  \end{itemize}
\end{block}

\end{frame}

%%%%%%%%%%%%%%%%%%%%%%%%%%%%%%%%%%%%%%%%%%%%%%%%%%%

\begin{frame}{Administrationspanel - Quick Actions \& Teknisk Bas}
\framesubtitle{Snabbåtkomst och systemarkitektur}

\begin{block}{4. Quick Actions Panel}
  \begin{itemize}
    \item \textbf{6 stora ikonknappar:}
      \begin{itemize}
        \item Skapa kund: Navigerar till onboarding-flöde
        \item Kör riskscan: Triggar bakgrundsjobb för all bokföringsdata
        \item Generera rapport: Exporterar PDF-sammanställning av alla tester
        \item Exportera data: CSV-nedladdning av aggregerad metadata
        \item Systeminställningar: Öppnar konfigurationsmodal
        \item Dokumentation: Länk till VALIDATION\_TESTS\_NEW.md och API-docs
      \end{itemize}
    \item \textbf{Design:} Färgglada ikoner, hover-effekter, tooltips med beskrivningar
  \end{itemize}
\end{block}

\begin{block}{Teknisk Implementation}
  \begin{itemize}
    \item \textbf{Frontend:} React + Tailwind CSS, routing via React Router (/admin)
    \item \textbf{Mock-data:} Hårdkodade JSON-objekt för demo (statistics, tickets, invoices, activities)
    \item \textbf{API-integration (framtida):}
      \begin{itemize}
        \item GET /api/admin/stats - Hämta dashboard-statistik
        \item GET /api/admin/support/tickets - Hämta supportärenden
        \item POST /api/admin/support/tickets/:id/resolve - Markera löst
        \item GET /api/admin/invoices - Hämta fakturaöversikt
        \item POST /api/admin/invoices/:id/remind - Skicka påminnelse
        \item GET /api/admin/activity - Hämta recent activity timeline
      \end{itemize}
    \item \textbf{Responsive design:} Fungerar på desktop (primär) och tablet
    \item \textbf{Åtkomstkontroll:} Endast admin-användare (role=admin) får tillgång
  \end{itemize}
\end{block}

\begin{block}{Mock-data för LIA-Demo}
  \begin{itemize}
    \item 3 mock användare (Anna, Erik, Lisa från SettingsPage)
    \item 5 supportärenden (2 open, 2 in-progress, 1 resolved)
    \item 10 fakturor (5 betalda, 3 väntande, 2 förfallna)
    \item 20 recent activities (login, upload, test flagged, support created)
    \item Statistik: 1 234 API calls idag, 45 678 denna månad, 12 aktiva sessioner, 87 totalt kunder
  \end{itemize}
\end{block}

\end{frame}



\end{document}



